% ===========================================================================
%  Chapitre 10 — Barycentre
%  PE 2026, composante GÉOMÉTRIE — série OSE
% ===========================================================================
\bmtchapitre{Barycentre}
  {Choisir une méthode pour résoudre un problème dans le plan à l'aide du
   barycentre de $n$ points pondérés, des fonctions vectorielle et scalaire de
   Leibniz et des lignes de niveau.}
  {Géométrie \textbullet\ environ 8 heures}

Le barycentre est né de la physique : c'est le point d'équilibre d'un système
de masses. Un fléau de balance chargé de deux poids bascule autour d'un point
unique, et ce point se calcule. Les mathématiques ont retenu l'idée en
supprimant la contrainte de positivité — les coefficients peuvent être
négatifs — et en ont fait un outil de géométrie.

Sa force tient à une seule propriété : une somme vectorielle pondérée se
réduit à un unique vecteur dès qu'on connaît le barycentre. Toutes les
configurations du problème de géométrie du baccalauréat en découlent, et
notamment les ensembles de points définis par une relation entre distances au
carré, qui reviennent chaque année.

% ---------------------------------------------------------------------------
\begin{revision}
Géométrie vectorielle de seconde et de première.

\begin{enumerate}[label=\textbf{\arabic*.}, leftmargin=*, itemsep=0.35em]
  \item Rappeler la relation de Chasles.
  \item $I$ étant le milieu de $[AB]$, exprimer $\vect{MA}+\vect{MB}$ en
        fonction de $\vect{MI}$.
  \item Que vaut $\vect{GA}+\vect{GB}+\vect{GC}$ lorsque $G$ est le centre de
        gravité du triangle $ABC$ ?
  \item Rappeler la définition du produit scalaire et l'identité
        $\norme{\vecu+\vecv}^{2} = \norme{\vecu}^{2}+2\scal{\vecu}{\vecv}
        +\norme{\vecv}^{2}$.
  \item Quel est l'ensemble des points $M$ tels que $MA = r$, avec $r>0$ ?
\end{enumerate}
\end{revision}

\begin{automatismes}
Sans calculatrice, sans brouillon.

\begin{enumerate}[label=\textbf{\alph*)}, leftmargin=*, itemsep=0.25em]
  \item $\vect{AB}+\vect{BC}$
  \item $\vect{MA}+\vect{MB}$ avec $I$ milieu de $[AB]$
  \item Coordonnées du milieu de $A(1\,;3)$ et $B(5\,;-1)$
  \item Affixe du milieu de $z_{A} = 2i$ et $z_{B} = 4$
  \item $\vect{GA}+\vect{GB}+\vect{GC}$, $G$ centre de gravité
  \item $\scal{\vecu}{\vecu}$
  \item $\norme{\vect{AB}}$ pour $A(0\,;0)$ et $B(3\,;4)$
  \item Ensemble des $M$ tels que $MA = 0$
\end{enumerate}
\end{automatismes}

\begin{corrige}[automatismes]
a) $\vect{AC}$ \quad b) $2\vect{MI}$ \quad c) $(3\,;1)$ \quad
d) $2+i$ \quad e) $\vecu$ nul \quad f) $\norme{\vecu}^{2}$ \quad
g) $5$ \quad h) le seul point $A$.
\end{corrige}

% ===========================================================================
\section{Barycentre de plusieurs points pondérés}

\begin{definition}[point pondéré, barycentre]
Un \textbf{point pondéré} est un couple $(A\,;\alpha)$ formé d'un point et
d'un réel appelé \textbf{coefficient}. Soit
$\left\{(A_{1}\,;\alpha_{1}),\dots,(A_{n}\,;\alpha_{n})\right\}$ un système de
points pondérés dont la somme des coefficients
$\alpha = \alpha_{1}+\cdots+\alpha_{n}$ est non nulle. Il existe alors un
unique point $G$, appelé \textbf{barycentre} du système, tel que
\[
  \alpha_{1}\vect{GA_{1}}+\alpha_{2}\vect{GA_{2}}+\cdots
  +\alpha_{n}\vect{GA_{n}} = \vecu_{0},
\]
où $\vecu_{0}$ désigne le vecteur nul.
\end{definition}

\begin{demonstration}
Fixons un point $O$. La relation de Chasles donne
$\sum \alpha_{i}\vect{GA_{i}} = \sum \alpha_{i}
\left(\vect{GO}+\vect{OA_{i}}\right)
= \alpha\,\vect{GO}+\sum \alpha_{i}\vect{OA_{i}}$. Cette somme est nulle si et
seulement si
\[
  \vect{OG} = \frac{1}{\alpha}\sum \alpha_{i}\vect{OA_{i}},
\]
ce qui détermine $G$ de façon unique.
\end{demonstration}

\begin{avertissement}
Si la somme des coefficients est nulle, il n'y a \emph{pas} de barycentre.
Cette situation n'est pas un accident : elle est exploitée à la section
suivante, car la somme vectorielle devient alors un vecteur constant. Vérifier
le signe de la somme des coefficients est donc le tout premier geste.
\end{avertissement}

\begin{propriete}[propriétés fondamentales]
\begin{enumerate}[label=\textbf{\arabic*.}, leftmargin=*, itemsep=0.15em]
  \item \emph{Homogénéité.} Multiplier tous les coefficients par un même réel
        non nul ne change pas le barycentre.
  \item \emph{Associativité.} On peut remplacer plusieurs points pondérés par
        leur barycentre partiel affecté de la somme de leurs coefficients.
  \item Si tous les coefficients sont égaux, $G$ est l'\textbf{isobarycentre} :
        milieu pour deux points, centre de gravité pour trois.
\end{enumerate}
\end{propriete}

\begin{visualisation}[le point d'équilibre]
\begin{center}
\begin{tikzpicture}[x=1.15cm, y=1.15cm, font=\footnotesize]
  \coordinate (A) at (0,0);
  \coordinate (B) at (5,0);
  \draw[bmtGris, line width=1pt] (A) -- (B);
  \filldraw[bmtIndigo] (A) circle (2.1pt) node[anchor=north] {$A$};
  \filldraw[bmtIndigo] (B) circle (2.1pt) node[anchor=north] {$B$};
  \filldraw[bmtLaterite] (2,0) circle (2.4pt);
  \node[bmtLaterite, anchor=north] at (2,-0.12) {$G$};
  \node[bmtIndigo, anchor=south] at (0,0.15) {$\alpha = 3$};
  \node[bmtIndigo, anchor=south] at (5,0.15) {$\beta = 2$};
  \draw[bmtLaterite, decorate, decoration={brace, amplitude=4pt}]
    (0,-0.55) -- (2,-0.55) node[midway, anchor=north, yshift=-4pt] {$2$};
  \draw[bmtLaterite, decorate, decoration={brace, amplitude=4pt}]
    (2,-0.55) -- (5,-0.55) node[midway, anchor=north, yshift=-4pt] {$3$};
  \draw[bmtVetiver, line width=1.4pt] (1.6,-1.35) -- (2.4,-1.35);
  \draw[bmtVetiver] (2,-1.35) -- (2,-1.1);
\end{tikzpicture}
\end{center}

Avec des coefficients $3$ et $2$, le barycentre partage le segment dans le
rapport \emph{inverse} des coefficients : ici $AG = \frac25\,AB$ et
$GB = \frac35\,AB$. Le point d'équilibre est donc plus proche du point le plus
lourd — c'est exactement la balance romaine, et le calcul confirme l'intuition
physique. Retenez la règle : plus un point pèse, plus le barycentre s'en
rapproche.
\end{visualisation}

\begin{entrainement}[construire un barycentre]
\exo[\niveau{1}] Construire le barycentre de
$\left\{(A\,;1),(B\,;3)\right\}$.

\exo[\niveau{1}] Construire le barycentre de
$\left\{(A\,;2),(B\,;-1)\right\}$ et le situer par rapport au segment.

\exo[\niveau{2}] $ABC$ étant un triangle, construire le barycentre de
$\left\{(A\,;1),(B\,;1),(C\,;2)\right\}$ en utilisant l'associativité.

\exo[\niveau{2}] Montrer que le barycentre de
$\left\{(A\,;1),(B\,;-1),(C\,;1)\right\}$ est le point $D$ tel que
$\vect{AD} = \vect{BC}$.

\exo[\niveau{3}] $ABC$ est un triangle et $F$ le barycentre de
$\left\{(A\,;2),(B\,;-1),(C\,;1)\right\}$. Soit $G$ le milieu de $[BC]$.
Montrer que $ABGF$ est un parallélogramme.
\end{entrainement}

% ===========================================================================
\section{Coordonnées et affixe du barycentre}

\begin{propriete}[expression analytique]
Dans un repère, les coordonnées du barycentre sont les moyennes pondérées des
coordonnées :
\[
  x_{G} = \frac{\sum \alpha_{i}x_{i}}{\sum \alpha_{i}},
  \qquad
  y_{G} = \frac{\sum \alpha_{i}y_{i}}{\sum \alpha_{i}} .
\]
Dans le plan complexe, l'affixe suit la même règle :
$z_{G} = \dfrac{\sum \alpha_{i}z_{i}}{\sum \alpha_{i}}$.
\end{propriete}

\begin{exemple}[une affixe de barycentre]
Le plan est rapporté au repère complexe
$\left(O\,;\vect{OA},\vect{OB}\right)$, de sorte que $z_{A} = 1$ et
$z_{B} = i$. Le barycentre $G$ de
$\left\{(O\,;4),(A\,;-1),(B\,;-1)\right\}$ a pour somme de coefficients
$4-1-1 = 2$, donc
\[
  z_{G} = \frac{4 \times 0-1 \times 1-1 \times i}{2} = -\frac{1+i}{2}.
\]
C'est le symétrique par rapport à $O$ du milieu de $[AB]$.
\end{exemple}

% ===========================================================================
\section{Réduction de la somme vectorielle}

\begin{theoreme}[fonction vectorielle de Leibniz]
Soit $f$ l'application qui à tout point $M$ associe
$\displaystyle\sum_{i=1}^{n}\alpha_{i}\vect{MA_{i}}$, et soit
$\alpha = \sum \alpha_{i}$.
\begin{itemize}[leftmargin=*, itemsep=0.1em]
  \item Si $\alpha \neq 0$ et si $G$ est le barycentre du système, alors
        $f(M) = \alpha\,\vect{MG}$ pour tout $M$.
  \item Si $\alpha = 0$, alors $f$ est \textbf{constante} : le vecteur
        $f(M)$ ne dépend pas de $M$.
\end{itemize}
\end{theoreme}

\begin{demonstration}
Chasles donne $\sum \alpha_{i}\vect{MA_{i}}
= \sum \alpha_{i}\left(\vect{MG}+\vect{GA_{i}}\right)
= \alpha\vect{MG}+\underbrace{\sum \alpha_{i}\vect{GA_{i}}}_{= \vecu_{0}}$
lorsque $G$ existe. Si $\alpha = 0$, on choisit un point fixe $O$ et l'on
écrit $\sum \alpha_{i}\vect{MA_{i}}
= \underbrace{\alpha}_{=0}\vect{MO}+\sum \alpha_{i}\vect{OA_{i}}$, quantité
indépendante de $M$.
\end{demonstration}

\begin{methode}{réduire une somme vectorielle}
Additionner les coefficients. Si la somme est non nulle, introduire le
barycentre et remplacer la somme par $\alpha\vect{MG}$. Si elle est nulle,
choisir l'un des points comme origine et développer par Chasles : le résultat
est un vecteur fixe, à exprimer à l'aide des données de la figure.
\end{methode}

\begin{exemple}[une égalité de normes]
Soit $ABCD$ un carré direct de côté $4$. Cherchons l'ensemble $(\Gamma)$ des
points $M$ tels que
\[
  \norme{-\vect{MA}+\vect{MB}+\vect{MD}}
  = \norme{-2\vect{MA}+\vect{MB}+\vect{MD}} .
\]
À gauche, la somme des coefficients vaut $1$ ; le barycentre du système
$\left\{(A\,;-1),(B\,;1),(D\,;1)\right\}$ est le point $C$, car
$\vect{AC} = \vect{AB}+\vect{AD}$. Le membre de gauche vaut donc $MC$.

À droite, la somme vaut $0$ : le vecteur est constant et l'on développe
\[
  -2\vect{MA}+\vect{MB}+\vect{MD}
  = \left(\vect{MB}-\vect{MA}\right)+\left(\vect{MD}-\vect{MA}\right)
  = \vect{AB}+\vect{AD} = \vect{AC},
\]
de norme $AC = 4\sqrt2$.

L'ensemble cherché est donc le cercle de centre $C$ et de rayon $4\sqrt2$.
\end{exemple}

% ===========================================================================
\section{Fonction scalaire de Leibniz}

\begin{theoreme}[réduction d'une somme de carrés]
Soit $g(M) = \displaystyle\sum_{i=1}^{n}\alpha_{i}\,MA_{i}^{2}$ et
$\alpha = \sum \alpha_{i}$.
\begin{itemize}[leftmargin=*, itemsep=0.1em]
  \item Si $\alpha \neq 0$, de barycentre $G$ :
        $g(M) = \alpha\,MG^{2}+g(G)$.
  \item Si $\alpha = 0$ : $g(M) = 2\,\scal{\vecu_{0}}{\vect{MO}}+g(O)$, où
        $\vecu_{0}$ est le vecteur constant de la section précédente et $O$ un
        point fixe quelconque.
\end{itemize}
\end{theoreme}

\begin{demonstration}
Dans le premier cas,
$MA_{i}^{2} = \norme{\vect{MG}+\vect{GA_{i}}}^{2}
= MG^{2}+2\scal{\vect{MG}}{\vect{GA_{i}}}+GA_{i}^{2}$. En multipliant par
$\alpha_{i}$ et en sommant, le terme central fait apparaître
$\sum \alpha_{i}\vect{GA_{i}}$, qui est nul. Il reste
$\alpha\,MG^{2}+\sum \alpha_{i}GA_{i}^{2} = \alpha MG^{2}+g(G)$.
\end{demonstration}

\begin{methode}{déterminer une ligne de niveau $\sum \alpha_{i}MA_{i}^{2} = k$}
Quatre étapes, dans cet ordre.
\begin{enumerate}[label=\textbf{\arabic*.}, leftmargin=*, itemsep=0.15em]
  \item Additionner les coefficients. Si la somme est nulle, la relation
        devient affine et l'ensemble est une droite.
  \item Sinon, construire le barycentre $G$.
  \item Calculer $g(G) = \sum \alpha_{i}GA_{i}^{2}$ à partir des données de la
        figure.
  \item Écrire $\alpha\,MG^{2} = k-g(G)$ et discuter suivant le signe du
        second membre : ensemble vide, point unique, ou cercle.
\end{enumerate}
\end{methode}

% ===========================================================================
\section{Lignes de niveau}

\begin{propriete}[les trois lignes de niveau du programme]
\begin{enumerate}[label=\textbf{\arabic*.}, leftmargin=*, itemsep=0.15em]
  \item $\left\{M \mid \sum \alpha_{i}MA_{i}^{2} = k\right\}$ avec
        $\alpha \neq 0$ : cercle de centre $G$, éventuellement réduit à un
        point ou vide.
  \item $\left\{M \mid \scal{\vect{AM}}{\vecu} = k\right\}$ : droite
        perpendiculaire à $\vecu$.
  \item $\left\{M \mid MA = \lambda MB\right\}$ avec $\lambda>0$ et
        $\lambda \neq 1$ : cercle, dont un diamètre est porté par $(AB)$.
\end{enumerate}
\end{propriete}

\begin{exemple}[une discussion complète]
$ABC$ est rectangle et isocèle en $A$ avec $AB = AC = 3$. Soit $D$ le
barycentre de $\left\{(A\,;1),(B\,;-1),(C\,;1)\right\}$ et, pour $k$ réel,
$(E_{k})$ l'ensemble des points $M$ tels que
$MA^{2}-MB^{2}+MC^{2} = k$.

La somme des coefficients vaut $1$, donc
$MA^{2}-MB^{2}+MC^{2} = MD^{2}+g(D)$ avec
$g(D) = DA^{2}-DB^{2}+DC^{2}$.

En repérant $A(0\,;0)$, $B(3\,;0)$ et $C(0\,;3)$, on trouve $D(-3\,;3)$, puis
$DA^{2} = 18$, $DB^{2} = 45$, $DC^{2} = 9$ : donc $g(D) = -18$ et
\[
  (E_{k}) : \quad MD^{2} = k+18 .
\]
Si $k<-18$, l'ensemble est vide ; si $k = -18$, il se réduit au point $D$ ;
si $k>-18$, c'est le cercle de centre $D$ et de rayon $\sqrt{k+18}$.
\end{exemple}

\begin{entrainement}[lignes de niveau]
\exo[\niveau{1}] $A$ et $B$ sont distants de $6$. Déterminer l'ensemble des
$M$ tels que $MA^{2}+MB^{2} = 26$.

\exo[\niveau{2}] Même triangle rectangle isocèle qu'à l'exemple. Déterminer
l'ensemble des $M$ tels que $MA^{2}+MB^{2}+MC^{2} = 27$.

\exo[\niveau{2}] Déterminer l'ensemble des $M$ tels que
$MA^{2}-MB^{2} = 12$, avec $AB = 4$.

\exo[\niveau{3}] $A$ et $B$ étant distants de $6$, déterminer l'ensemble des
$M$ tels que $MA = 2MB$.
\end{entrainement}

% ===========================================================================
\begin{annales}
Le barycentre ouvre le problème de géométrie, presque toujours par la
construction d'un point puis par une ligne de niveau. Les énoncés qui suivent
sont repris de sujets réellement posés ; leur provenance est indiquée à chaque
fois.

\exoann{Baccalauréat, Madagascar, série S, session 2021 — problème 1, partie A}
$ABC$ est un triangle rectangle en $A$ tel que
$\left(\vect{AB},\vect{AC}\right) = \frac{\pi}{2}$ et $AB = AC = 3$~cm.
\begin{enumerate}[label=\textbf{\alph*)}, leftmargin=*, itemsep=0.15em]
  \item Déterminer et construire le point $D$, barycentre du système
        $\left\{(A\,;1),(B\,;-1),(C\,;1)\right\}$.
  \item Discuter, suivant les valeurs du paramètre réel $k$, la nature de
        l'ensemble $(E_{k})$ des points $M$ du plan vérifiant
        $MA^{2}-MB^{2}+MC^{2} = k$. \emph{On pourra montrer que
        $MD^{2} = k+18$.}
\end{enumerate}

\exoann{Baccalauréat, Madagascar, série S, session 2021 — problème 1, partie A, suite}
Avec les mêmes données, $E$ désigne le symétrique de $B$ par rapport à $A$.
Soit $(\mathcal C)$ le cercle de centre $C$ et de rayon $CB$, et $F$ son
second point d'intersection avec la droite $(BC)$. Montrer que
$\vect{ED} = \frac12\vect{EF}$.

\exoann{Baccalauréat, Madagascar, série S, session 2023 — problème 1, partie I}
$ABC$ est rectangle et isocèle en $A$, $AB = 4$~cm,
$\left(\vect{AB},\vect{AC}\right) = \frac{\pi}{2}$.
\begin{enumerate}[label=\textbf{\alph*)}, leftmargin=*, itemsep=0.15em]
  \item Déterminer les points
        $D = \bary{(A\,;1),(B\,;-1),(C\,;-1)}$ et
        $G = \bary{(B\,;1-\sqrt2),(A\,;\sqrt2)}$.
  \item Déterminer et construire l'ensemble $(\Gamma)$ des points $M$ du plan
        vérifiant $MA^{2}-MB^{2}-MC^{2} = -12$.
\end{enumerate}

\exoann{Baccalauréat, Madagascar, série S, session 2026 — problème 1, partie I}
$ABC$ est un triangle direct rectangle en $A$ tel que $AC = 4$~cm et
$BC = 8$~cm.
\begin{enumerate}[label=\textbf{\alph*)}, leftmargin=*, itemsep=0.15em]
  \item Déterminer et construire $G$, barycentre du système
        $\left\{(A\,;-1),(B\,;1),(C\,;1)\right\}$.
  \item Déterminer l'ensemble $(\mathcal C)$ des points $M$ du plan tels que
        $-MA^{2}+MB^{2}+MC^{2} = 16$.
\end{enumerate}

\exoann{Baccalauréat, Madagascar, série S, session 2026 — problème 1, partie I, suite}
Avec les mêmes données, $O$ désigne le milieu de $[AC]$. Déterminer et
construire le point $\Omega$, barycentre de
$\left\{(A\,;2),(B\,;-1)\right\}$.

\exoann{Baccalauréat, Madagascar, série S, session 2022 — problème 1, partie A}
$ABCD$ est un carré direct de côté $AB = 4$~cm.
\begin{enumerate}[label=\textbf{\alph*)}, leftmargin=*, itemsep=0.15em]
  \item Prouver que $C$ est le barycentre du système
        $\left\{(A\,;-1),(B\,;1),(D\,;1)\right\}$.
  \item Déterminer et construire l'ensemble $(\Gamma)$ des points $M$ du plan
        vérifiant
        $\norme{-\vect{MA}+\vect{MB}+\vect{MD}}
        = \norme{-2\vect{MA}+\vect{MB}+\vect{MD}}$.
\end{enumerate}

\exoann{Baccalauréat blanc, lycée Philibert Tsiranana, Terminale S, 2021-2022 — problème 1, partie A}
Dans le plan orienté, $OAB$ est un triangle tel que $OA = OB = a$, $a>0$, et
$\left(\vect{OA},\vect{OB}\right) = \frac{\pi}{2}$.
\begin{enumerate}[label=\textbf{\alph*)}, leftmargin=*, itemsep=0.15em]
  \item Déterminer et construire le barycentre $G$ du système
        $\left\{(O\,;4),(A\,;-1),(B\,;-1)\right\}$.
  \item Déterminer et construire l'ensemble $(E)$ des points $M$ du plan
        vérifiant $4MO^{2}-MA^{2}-MB^{2} = 2a^{2}$.
\end{enumerate}

\exoann{Baccalauréat, Congo-Brazzaville, série C, session 2022 — exercice 2}
$ABC$ est un triangle équilatéral direct de centre de gravité $G$, $O$ est le
milieu de $[AB]$, $I$ le symétrique de $C$ par rapport à $O$, et $BC = 5$~cm.
Soit $(\mathcal C)$ le cercle de diamètre $[IG]$. Montrer que
$(\mathcal C)$ est l'ensemble des points $M$ du plan tels que
$MC = 2MO$.

\exoann{Baccalauréat, Congo-Brazzaville, série C, session 2022 — exercice 2, suite}
Avec les mêmes données, soit $h$ l'application qui à tout point $M$ associe le
point $M'$ tel que
$\vect{MM'} = \vect{MA}+\vect{MB}+\vect{MC}$.
\begin{enumerate}[label=\textbf{\alph*)}, leftmargin=*, itemsep=0.15em]
  \item Démontrer que $\vect{GM'} = -2\vect{GM}$.
  \item En déduire la nature de $h$.
\end{enumerate}

\exoann{Baccalauréat, Madagascar, série C, session 2012 — problème 1, partie I}
$ABC$ est un triangle rectangle et isocèle en $A$ tel que $AB = AC = 2$. On
note $G$ le milieu de $[BC]$. Déterminer et construire le barycentre $F$ des
points $A$, $B$ et $C$ affectés respectivement des coefficients $2$, $-1$ et
$1$, puis montrer que $ABGF$ est un parallélogramme.
\end{annales}

\begin{corrige}[annales]
\textbf{1. a)} La somme des coefficients vaut $1$, donc
$\vect{AD} = -\vect{AB}+\vect{AC} = \vect{BC}$ : $D$ est l'image de $A$ par
la translation de vecteur $\vect{BC}$, et $ABCD$ est un parallélogramme.
\textbf{b)} En repérant $A(0\,;0)$, $B(3\,;0)$, $C(0\,;3)$, on a
$D(-3\,;3)$, puis $DA^{2} = 18$, $DB^{2} = 45$ et $DC^{2} = 9$, d'où
$g(D) = -18$. Ainsi $(E_{k}) : MD^{2} = k+18$ : ensemble vide si $k<-18$,
point $D$ si $k = -18$, cercle de centre $D$ et de rayon $\sqrt{k+18}$ si
$k>-18$.

\textbf{2.} Avec le même repère, $E(-3\,;0)$ et $D(-3\,;3)$, donc
$\vect{ED}$ a pour coordonnées $(0\,;3)$. Le point $F$, second point
d'intersection de $(BC)$ avec le cercle de centre $C$ et de rayon $CB$, est le
symétrique de $B$ par rapport à $C$, soit $F(-3\,;6)$ ; alors $\vect{EF}$ a
pour coordonnées $(0\,;6)$. On conclut
$\vect{ED} = \frac12\vect{EF}$.

\textbf{3. a)} La somme des coefficients de $D$ vaut $-1$, donc
$\vect{AD} = \vect{AB}+\vect{AC}$ : $D$ est le quatrième sommet du carré
construit sur $[AB]$ et $[AC]$, de coordonnées $(4\,;4)$. Pour $G$, la somme
vaut $1$ et $\vect{AG} = (1-\sqrt2)\vect{AB}$ : $G$ est sur $(AB)$, d'abscisse
$4-4\sqrt2$.
\textbf{b)} La somme des coefficients de $(\Gamma)$ vaut $-1$ et le barycentre
est $D$. Or $DA^{2} = 32$, $DB^{2} = 16$ et $DC^{2} = 16$, donc
$g(D) = 32-16-16 = 0$ et $(\Gamma) : -MD^{2} = -12$, soit $MD = 2\sqrt3$ :
c'est le cercle de centre $D$ et de rayon $2\sqrt3$.

\textbf{4. a)} La somme vaut $1$ et $\vect{AG} = \vect{AB}+\vect{AC}$ : $G$
est le quatrième sommet du rectangle construit sur $[AB]$ et $[AC]$. Avec
$A(0\,;0)$, $C(0\,;4)$ et $B\left(4\sqrt3\,;0\right)$ — car
$AB = \sqrt{64-16} = 4\sqrt3$ —, on obtient $G\left(4\sqrt3\,;4\right)$.
\textbf{b)} $GA^{2} = 64$, $GB^{2} = 16$, $GC^{2} = 48$, donc
$g(G) = -64+16+48 = 0$ et l'ensemble est défini par $MG^{2} = 16$ : c'est le
cercle de centre $G$ et de rayon $4$.

\textbf{5.} La somme vaut $1$, donc $\vect{A\Omega} = -\vect{AB}$ : $\Omega$
est le symétrique de $B$ par rapport à $A$, de coordonnées
$\left(-4\sqrt3\,;0\right)$.

\textbf{6. a)} La somme des coefficients vaut $1$ et
$\vect{AC} = \vect{AB}+\vect{AD}$ dans le carré : $C$ est bien le barycentre
annoncé.
\textbf{b)} Le membre de gauche vaut $MC$ d'après la question précédente. À
droite, la somme des coefficients est nulle, donc le vecteur est constant et
vaut $\vect{AB}+\vect{AD} = \vect{AC}$, de norme $4\sqrt2$. L'ensemble
$(\Gamma)$ est donc le cercle de centre $C$ et de rayon $4\sqrt2$.

\textbf{7. a)} La somme vaut $2$ et
$\vect{OG} = -\frac12\left(\vect{OA}+\vect{OB}\right) = -\vect{OI}$, où $I$
est le milieu de $[AB]$ : $G$ est le symétrique de $I$ par rapport à $O$.
\textbf{b)} En repérant $O(0\,;0)$, $A(a\,;0)$, $B(0\,;a)$, on trouve
$G\left(-\frac a2\,;-\frac a2\right)$, puis $GO^{2} = \frac{a^{2}}{2}$,
$GA^{2} = GB^{2} = \frac{5a^{2}}{2}$, d'où $g(G) = -3a^{2}$. La relation
devient $2MG^{2} = 2a^{2}+3a^{2} = 5a^{2}$, soit
$MG = \frac{a\sqrt{10}}{2}$ : c'est le cercle de centre $G$ et de ce rayon.

\textbf{8.} $MC = 2MO$ équivaut à $MC^{2}-4MO^{2} = 0$. La somme des
coefficients vaut $-3$ ; en plaçant $O$ à l'origine et $C$ sur un axe, les
points de la droite $(OC)$ vérifiant la relation sont d'abscisses $-c$ et
$\frac c3$, c'est-à-dire exactement $I$ et $G$. L'ensemble est donc le cercle
de diamètre $[IG]$.

\textbf{9. a)} La somme des coefficients de $\vect{MA}+\vect{MB}+\vect{MC}$
vaut $3$, donc cette somme égale $3\vect{MG}$. Ainsi
$\vect{GM'} = \vect{GM}+\vect{MM'} = \vect{GM}+3\vect{MG} = -2\vect{GM}$.
\textbf{b)} L'application $h$ est donc l'homothétie de centre $G$ et de
rapport $-2$.

\textbf{10.} La somme vaut $2$, donc
$\vect{AF} = \frac12\left(-\vect{AB}+\vect{AC}\right) = \frac12\vect{BC}$. Or
$\vect{BG} = \frac12\vect{BC}$ puisque $G$ est le milieu de $[BC]$. Les
vecteurs $\vect{AF}$ et $\vect{BG}$ sont égaux : $ABGF$ est un
parallélogramme.
\end{corrige}

% ===========================================================================
\begin{jecherche}[le point de collecte du riz]
Trois villages $A$, $B$ et $C$ sont situés aux sommets d'un triangle. Chaque
année, ils apportent respectivement $2$, $3$ et $5$ tonnes de riz à un point
de collecte $P$, et le coût du transport est proportionnel au produit de la
masse par le carré de la distance parcourue. Le coût total s'écrit donc
\[
  f(P) = 2\,PA^{2}+3\,PB^{2}+5\,PC^{2},
\]
à une constante près.

\begin{enumerate}[label=\textbf{\arabic*.}, leftmargin=*, itemsep=0.3em]
  \item Justifier l'existence du barycentre $G$ de
        $\left\{(A\,;2),(B\,;3),(C\,;5)\right\}$ et le construire.
  \item Exprimer $f(P)$ en fonction de $PG$ et de $f(G)$.
  \item En déduire la position du point de collecte qui minimise le coût.
  \item On place $A(0\,;0)$, $B(6\,;0)$ et $C(2\,;8)$. Calculer les
        coordonnées de $G$ puis la valeur minimale de $f$.
  \item Les villages acceptent un surcoût de $40$ unités par rapport au
        minimum. Décrire l'ensemble des emplacements possibles.
  \item Le village $C$ double sa production. Sans refaire tout le calcul, dire
        dans quelle direction le point de collecte se déplace.
\end{enumerate}
\end{jecherche}

\begin{corrige}[le point de collecte du riz]
\textbf{1.} La somme des coefficients vaut $10 \neq 0$ : le barycentre existe.
On le construit par associativité, en remplaçant $A$ et $B$ par leur
barycentre partiel affecté de $5$.

\textbf{2.} La fonction scalaire de Leibniz donne
$f(P) = 10\,PG^{2}+f(G)$.

\textbf{3.} Le terme $10PG^{2}$ est positif ou nul, nul seulement en $P = G$ :
le coût est minimal lorsque le point de collecte est placé en $G$, et le
minimum vaut $f(G)$.

\textbf{4.} $x_{G} = \frac{2 \times 0+3 \times 6+5 \times 2}{10} = 2{,}8$ et
$y_{G} = \frac{0+0+5 \times 8}{10} = 4$. Alors
$GA^{2} = 2{,}8^{2}+4^{2} = 23{,}84$,
$GB^{2} = 3{,}2^{2}+4^{2} = 26{,}24$,
$GC^{2} = 0{,}8^{2}+4^{2} = 16{,}64$, d'où
$f(G) = 2(23{,}84)+3(26{,}24)+5(16{,}64) = 209{,}6$.

\textbf{5.} $f(P) \leq f(G)+40$ équivaut à $10PG^{2} \leq 40$, soit
$PG \leq 2$ : les emplacements acceptables forment le disque de centre $G$ et
de rayon $2$.

\textbf{6.} Augmenter le coefficient de $C$ rapproche le barycentre de $C$ :
le point de collecte se déplace le long du segment $[GC]$, vers $C$. Le
calcul le confirmerait, mais la propriété d'équilibre suffit à le prévoir.
\end{corrige}

% ===========================================================================
\begin{jemeteste}
Répondre par vrai ou faux, en justifiant en une ligne.

\begin{enumerate}[label=\textbf{\arabic*.}, leftmargin=*, itemsep=0.28em]
  \item Tout système de points pondérés admet un barycentre.
  \item Le barycentre de $\left\{(A\,;1),(B\,;1)\right\}$ est le milieu de
        $[AB]$.
  \item Multiplier tous les coefficients par $-2$ change le barycentre.
  \item Si la somme des coefficients est nulle,
        $\sum \alpha_{i}\vect{MA_{i}}$ ne dépend pas de $M$.
  \item L'ensemble des $M$ tels que $MA^{2}+MB^{2} = k$ est toujours un
        cercle.
  \item Le barycentre de trois points affectés de coefficients positifs est
        intérieur au triangle.
\end{enumerate}
\end{jemeteste}

\begin{corrige}[je me teste]
\textbf{1.} Faux — il faut que la somme des coefficients soit non nulle.
\textbf{2.} Vrai.
\textbf{3.} Faux — le barycentre est inchangé par homogénéité.
\textbf{4.} Vrai — c'est la fonction vectorielle de Leibniz dans le cas
$\alpha = 0$.
\textbf{5.} Faux — il peut être vide ou réduit à un point selon $k$.
\textbf{6.} Vrai — il s'écrit comme moyenne pondérée des trois sommets.
\end{corrige}

% ===========================================================================
\begin{commeaubac}
\bmtsource{Baccalauréat, série S, session 2026 — problème 1, partie I}
$ABC$ est un triangle direct rectangle en $A$ tel que $AC = 4$~cm et
$BC = 8$~cm. On note $O$ le milieu de $[AC]$ et $D$ celui de $[BC]$.

\begin{enumerate}[label=\textbf{\arabic*.}, leftmargin=*, itemsep=0.25em]
  \item Montrer que le triangle $CAD$ est équilatéral.
        \hfill\textup{(0,5 pt)}
  \item Déterminer et construire le point $\Omega$, barycentre de
        $\left\{(A\,;2),(B\,;-1)\right\}$. \hfill\textup{(0,5 pt)}
  \item Déterminer et construire le point $G$, barycentre de
        $\left\{(A\,;-1),(B\,;1),(C\,;1)\right\}$. \hfill\textup{(0,5 pt)}
  \item Déterminer l'ensemble $(\mathcal C)$ des points $M$ du plan tels que
        $-MA^{2}+MB^{2}+MC^{2} = 16$. \hfill\textup{(0,75 pt)}
\end{enumerate}
\end{commeaubac}

\begin{corrige}[comme au baccalauréat]
\textbf{1.} Dans le triangle rectangle en $A$, la médiane issue de $A$ vaut la
moitié de l'hypoténuse : $AD = \frac{BC}{2} = 4$. Or $CD = \frac{BC}{2} = 4$
et $AC = 4$ : les trois côtés du triangle $CAD$ mesurent $4$~cm, il est
équilatéral.

\textbf{2.} La somme des coefficients vaut $1$, donc
$\vect{A\Omega} = -\vect{AB}$ : $\Omega$ est le symétrique de $B$ par rapport
à $A$.

\textbf{3.} La somme vaut $1$ et $\vect{AG} = \vect{AB}+\vect{AC}$ : $G$ est
le quatrième sommet du rectangle de côtés $[AB]$ et $[AC]$.

\textbf{4.} Avec $A(0\,;0)$, $C(0\,;4)$ et $B\left(4\sqrt3\,;0\right)$, on
obtient $G\left(4\sqrt3\,;4\right)$, puis $GA^{2} = 64$, $GB^{2} = 16$ et
$GC^{2} = 48$, donc $-GA^{2}+GB^{2}+GC^{2} = 0$. La relation devient
$MG^{2} = 16$ : $(\mathcal C)$ est le cercle de centre $G$ et de rayon $4$.
\end{corrige}

% ===========================================================================
\begin{notepourlaclasse}
Le premier obstacle n'est pas le barycentre lui-même mais le réflexe
d'additionner les coefficients avant tout. Faites-en une consigne écrite au
tableau, encadrée, et rappelée à chaque exercice pendant deux séances : la
somme est-elle nulle ou non ? Tout le reste en découle.

Le second obstacle est le calcul de $g(G)$ dans la fonction scalaire de
Leibniz. Beaucoup d'élèves cherchent à le mener vectoriellement et s'y
perdent ; le repère orthonormé bien choisi — un sommet à l'origine, un côté
sur un axe — ramène le calcul à trois distances. Imposez ce choix.

Enfin, la discussion suivant $k$ est systématiquement bâclée : les élèves
donnent le cercle et oublient les deux cas dégénérés. Comptez-les
explicitement dans votre barème, une fois, et l'oubli disparaît.
\end{notepourlaclasse}
